\section{Calibration of the Lotka--Volterra example}
\label{sec:calibration}
The source's Eq.~(22) defines
\begin{equation}
 L=\sqrt d\max_i\sup_{x\in\Omega}\norm{\nabla F_i(x)}_2,
 \qquad H=\sqrt d\max_i\sup_{x\in\Omega}\norm{D^2 F_i(x)}_F.
 \label{eq:source-LH}
\end{equation}
Its effective width and parameter, source Eqs.~(23), (25), are
\begin{equation}
 \sigma_F^2=\frac{\norm{\Sigma}_F}{\sqrt d},\qquad
 R=\frac{\sqrt d\,\sigma_F H}{2L}.
 \label{eq:source-R}
\end{equation}
For $\Sigma=\sigma^2 I$, the definition gives $\sigma_F=\sigma$.
The logarithmic Lotka--Volterra coordinates in the source give
\begin{equation}
 F_q(q,p)=-\alpha+\beta e^{-p},\qquad
 F_p(q,p)=\delta-\gamma e^{-q}.
 \label{eq:LV-field}
\end{equation}
The Euclidean norm of each row gradient equals the Frobenius norm of its
row Hessian at every point:
\begin{equation}
 \norm{\nabla F_q}_2=\norm{D^2F_q}_F=\beta e^{-p},\qquad
 \norm{\nabla F_p}_2=\norm{D^2F_p}_F=\gamma e^{-q}.
\end{equation}
Consequently $H=L$ on any common domain used in
Eq.~\eqref{eq:source-LH}, and for this two-dimensional example
\begin{equation}
 \frac R\sigma=\frac{1}{\sqrt2}=0.7071067812\ldots.
 \label{eq:correct-R}
\end{equation}
This differs from the source's Eq.~(45), which states
$R/\sigma=\gamma/\beta=2.29$. The inconsistency is with the displayed
definitions, independent of whether the common supremum is evaluated globally
or the same rowwise maxima are evaluated locally. With the source parameters
$\beta=e^{-1/3}$, $\gamma=e^{1/2}$, the ratio $\gamma/\beta$ is actually
$2.300975891\ldots$; that small decimal discrepancy is secondary to using the
wrong ratio. At $\sigma=0.01$, Eq.~\eqref{eq:correct-R} gives
$R=0.0070710678\ldots$, instead of $0.0229$.

Both values are small, so this calculation does not overturn the qualitative
success of that example. It changes the numerical calibration offered as
evidence for the proposed criterion. It also prevents interpreting the
reported value as a direct evaluation of the stated $L,H,\sigma_F$ definitions.

The paper's recommendation of roughly ten grid points across a packet width
is explicitly an order-of-magnitude numerical choice. We do not turn that
choice into a sharp theorem or claim a benchmark fails solely because a
chosen width spans fewer cells.
